% ==============================================================================
% BIBLIOGRAPHY ADDITIONS FOR "The Observable Universe as a Condensate
% Expanding into a Universal Protofluid"
%
% Compiled by: physics-agent
% Date: 2026-02-08
%
% These entries fill critical gaps identified in the doc-agent structural review
% and physics-agent content review. Each entry includes an annotation explaining
% its relevance to the paper's condensate cosmology framework.
%
% Format matches existing bibliography style in universal_condensate_draft_v01.tex
% ==============================================================================

% ==============================================================================
% SECTION 1: REQUIRED ENTRIES (per team review)
% ==============================================================================

% --- Entry 1: Volovik - The Universe in a Helium Droplet ---
% ANNOTATION: This is THE canonical reference for emergent spacetime from quantum
% liquids. Volovik systematically develops the connection between condensed matter
% physics (superfluid He-3, BEC) and gravitational dynamics, showing how effective
% metrics, gauge fields, and even fermions can emerge from condensate physics.
% This book provides the theoretical foundation for the paper's central claim that
% the universe is a condensate with emergent gravitational response. Mandatory
% citation for any condensate cosmology framework.
\bibitem{Volovik2003}
G.~E.~Volovik, \emph{The Universe in a Helium Droplet} (Oxford University Press, Oxford, 2003; paperback edition 2009). ISBN: 978-0-19-956484-2.

% --- Entry 2: Steinhauer - Hawking radiation in BEC ---
% ANNOTATION: This is the landmark experimental verification of analogue Hawking
% radiation in a Bose-Einstein condensate system. Steinhauer demonstrated that
% quantum correlations (entanglement) between Hawking particles and their partners
% arise at an acoustic horizon in a flowing BEC, confirming a key prediction of
% analogue gravity. This paper validates that BEC systems exhibit genuine quantum
% gravitational phenomena in the laboratory, strongly supporting the physical
% plausibility of condensate-based cosmology.
\bibitem{Steinhauer2016}
J.~Steinhauer, ``Observation of quantum Hawking radiation and its entanglement in an analogue black hole,'' \emph{Nature Physics}\ \textbf{12}, 959--965 (2016). arXiv:1510.00621 [gr-qc].

% --- Entry 3: Coleman and De Luccia - Vacuum decay with gravity ---
% ANNOTATION: This is the standard reference for first-order phase transitions
% (bubble nucleation) including gravitational back-reaction. The Coleman-De Luccia
% instanton describes quantum tunneling from a false vacuum to a true vacuum in
% the presence of gravity. This is directly relevant to the paper's interpretation
% of the Big Bang as a nucleation event---the protofluid-to-condensate transition
% is precisely this type of first-order phase transition. Essential for the
% theoretical grounding of Postulate 2 (Phase-Transition Origin of the Big Bang).
\bibitem{ColemanDeLuccia1980}
S.~Coleman and F.~De~Luccia, ``Gravitational effects on and of vacuum decay,'' \emph{Phys.\ Rev.\ D}\ \textbf{21}, 3305--3315 (1980).

% --- Entry 4: Garay et al. - BEC analogue black holes ---
% ANNOTATION: This foundational paper proposed using Bose-Einstein condensates
% to create sonic analogues of gravitational black holes. Garay, Anglin, Cirac,
% and Zoller showed that the Bogoliubov excitations of a flowing BEC propagate
% on an effective acoustic metric that can exhibit horizons and ergoregions. This
% established the theoretical basis for laboratory tests of analogue gravity
% phenomena, including the Steinhauer experiment cited above. Essential reference
% for the BEC-gravity connection underlying the paper's framework.
\bibitem{GarayEtAl2000}
L.~J.~Garay, J.~R.~Anglin, J.~I.~Cirac, and P.~Zoller, ``Sonic analogs of gravitational black holes in Bose-Einstein condensates,'' \emph{Phys.\ Rev.\ Lett.}\ \textbf{85}, 4643--4647 (2000). arXiv:gr-qc/0002015.

% --- Entry 5: Israel - Junction conditions ---
% ANNOTATION: This is the fundamental paper establishing the junction conditions
% (also called Israel junction conditions or Israel matching conditions) for
% singular hypersurfaces and thin shells in general relativity. These conditions
% govern how the metric and stress-energy tensor must match across a surface of
% discontinuity, which is exactly what the paper needs for the conversion boundary
% Sigma. The Rankine-Hugoniot relations in Section 3.1 are the cosmological
% application of Israel's formalism. Mandatory reference for rigorous treatment
% of the boundary dynamics.
\bibitem{Israel1966}
W.~Israel, ``Singular hypersurfaces and thin shells in general relativity,'' \emph{Nuovo Cimento B}\ \textbf{44}, 1--14 (1966); erratum \emph{ibid.}\ \textbf{48}, 463 (1967).

% --- Entry 6: Jacobson - Thermodynamics of spacetime ---
% ANNOTATION: Jacobson's seminal paper shows that the Einstein field equations
% can be derived as an equation of state from thermodynamic arguments applied
% to local Rindler horizons. This supports the paper's mechanical-geometric
% linkage (Postulate 3) by providing an independent derivation of Einstein gravity
% from thermodynamic/mechanical principles. If spacetime has thermodynamic degrees
% of freedom, it strengthens the case that gravity emerges from an underlying
% medium with mechanical properties (like a condensate).
\bibitem{Jacobson1995}
T.~Jacobson, ``Thermodynamics of spacetime: The Einstein equation of state,'' \emph{Phys.\ Rev.\ Lett.}\ \textbf{75}, 1260--1263 (1995). arXiv:gr-qc/9504004.

% --- Entry 7: Volovik - Fermi-point scenario ---
% ANNOTATION: This paper develops the Fermi-point mechanism for emergent relativity
% and gauge fields from topological defects in the Fermi surface of quantum liquids.
% Volovik shows that near Fermi points (zeros of the energy spectrum), low-energy
% excitations necessarily obey Lorentz-invariant dynamics with an emergent speed
% of light set by the Fermi velocity. This provides a microphysical mechanism for
% why the acoustic metric in a condensate becomes Lorentzian, addressing a key
% gap identified in the physics review (Section 5).
\bibitem{Volovik2008}
G.~E.~Volovik, ``Emergent physics: Fermi point scenario,'' \emph{Phil.\ Trans.\ R.\ Soc.\ A}\ \textbf{366}, 2935--2951 (2008). arXiv:0801.0724 [cond-mat.other].

% ==============================================================================
% SECTION 2: ADDITIONAL RECOMMENDED ENTRIES
% These fill remaining gaps in BEC cosmology, Gross-Pitaevskii hydrodynamics,
% relativistic shock physics, and recent EDE developments
% ==============================================================================

% --- BEC cosmology and superfluid universe models ---

% ANNOTATION: Unruh's original paper establishing that accelerated observers
% detect thermal radiation (Unruh effect). This is foundational for analogue
% gravity, as it connects acceleration to temperature, and by the equivalence
% principle, to gravitational phenomena. Already implicitly referenced in the
% draft via AnalogueGravityReview, but worth citing directly.
\bibitem{Unruh1976}
W.~G.~Unruh, ``Notes on black-hole evaporation,'' \emph{Phys.\ Rev.\ D}\ \textbf{14}, 870--892 (1976).

% ANNOTATION: Unruh's key paper establishing the acoustic metric formalism---
% that sound waves in a moving fluid propagate on an effective curved spacetime
% geometry. This is the theoretical origin of "dumb holes" (acoustic black holes)
% and the foundation for all subsequent analogue gravity work.
\bibitem{Unruh1981}
W.~G.~Unruh, ``Experimental Black-Hole Evaporation?'' \emph{Phys.\ Rev.\ Lett.}\ \textbf{46}, 1351--1353 (1981).

% ANNOTATION: Visser's pedagogical paper on acoustic black holes, providing
% a clear derivation of the acoustic metric from fluid dynamics and discussing
% the conditions for horizon formation. Complements the Barcelo et al. review
% with more focus on the acoustic metric derivation.
\bibitem{Visser1998}
M.~Visser, ``Acoustic black holes: horizons, ergospheres and Hawking radiation,'' \emph{Class.\ Quantum Grav.}\ \textbf{15}, 1767--1791 (1998). arXiv:gr-qc/9712010.

% --- Gross-Pitaevskii hydrodynamics and acoustic metrics ---

% ANNOTATION: Dalfovo et al.'s comprehensive review of BEC theory and experiment,
% including the Gross-Pitaevskii equation, collective excitations, and the
% hydrodynamic limit. Essential background for understanding the microphysics
% underlying the condensate framework.
\bibitem{DalfovoEtAl1999}
F.~Dalfovo, S.~Giorgini, L.~P.~Pitaevskii, and S.~Stringari, ``Theory of Bose-Einstein condensation in trapped gases,'' \emph{Rev.\ Mod.\ Phys.}\ \textbf{71}, 463--512 (1999). arXiv:cond-mat/9806038.

% ANNOTATION: Foundational paper on excitations in weakly interacting Bose gases,
% establishing the Bogoliubov theory that underlies the acoustic metric derivation.
% The Bogoliubov dispersion relation shows how phonons become the relevant
% low-energy excitations in a BEC.
\bibitem{Bogoliubov1947}
N.~N.~Bogoliubov, ``On the theory of superfluidity,'' \emph{J.\ Phys.\ (USSR)}\ \textbf{11}, 23 (1947). [Reprinted in various collections.]

% --- Relativistic Rankine-Hugoniot and shock physics ---

% ANNOTATION: Taub's paper generalizing shock-wave theory to relativistic fluids.
% The relativistic Rankine-Hugoniot (Taub) adiabat is essential for treating
% the conversion boundary in a fully relativistic framework. Directly relevant
% to Section 3.1 and Appendix A.
\bibitem{Taub1948}
A.~H.~Taub, ``Relativistic Rankine-Hugoniot Equations,'' \emph{Phys.\ Rev.}\ \textbf{74}, 328--334 (1948).

% ANNOTATION: Lichnerowicz's rigorous mathematical treatment of relativistic
% shock waves. Provides the formal framework for the generalized Rankine-Hugoniot
% relations with surface stress-energy (Equation 5 in the draft).
\bibitem{Lichnerowicz1967}
A.~Lichnerowicz, \emph{Relativistic Hydrodynamics and Magnetohydrodynamics} (W.~A.~Benjamin, New York, 1967).

% --- Additional EDE and Hubble tension references (2024-2025) ---

% ANNOTATION: Original EDE paper by Poulin et al. proposing an oscillating
% scalar field that becomes dynamically important before recombination. This
% established the parametric framework (f_EDE, z_c, theta_i) that the paper's
% u_mech phenomenology mirrors. Essential reference for the EDE-like mechanism.
\bibitem{Poulin2019}
V.~Poulin, T.~L.~Smith, T.~Karwal, and M.~Kamionkowski, ``Early Dark Energy can Resolve the Hubble Tension,'' \emph{Phys.\ Rev.\ Lett.}\ \textbf{122}, 221301 (2019). arXiv:1811.04083.

% ANNOTATION: Smith, Poulin, and Amin's paper on oscillating scalar fields
% (axion-like particles) and the Hubble tension. Provides more detailed
% modeling of the EDE dynamics and perturbation theory.
\bibitem{SmithPoulinAmin2020}
T.~L.~Smith, V.~Poulin, and M.~A.~Amin, ``Oscillating scalar fields and the Hubble tension: A resolution with novel signatures,'' \emph{Phys.\ Rev.\ D}\ \textbf{101}, 063523 (2020). arXiv:1908.06995.

% ANNOTATION: DESI Year 1 BAO results represent the most recent BAO measurements.
% Essential for up-to-date validation of the condensate framework against
% late-time cosmological observables.
\bibitem{DESI2024}
DESI Collaboration (A.~G.~Adame \emph{et al.}), ``DESI 2024 VI: Cosmological Constraints from the Measurements of Baryon Acoustic Oscillations,'' arXiv:2404.03002 (2024).

% ANNOTATION: Analysis of DESI+CMB data in the context of the Hubble tension,
% representing the state of the art in combined cosmological constraints.
\bibitem{DESICMB2024}
DESI Collaboration (A.~G.~Adame \emph{et al.}), ``DESI 2024 VII: Cosmological Constraints from Full-Shape Analyses of the Two-Point Clustering Statistics,'' arXiv:2411.12022 (2024).

% --- Phase transitions in the early universe ---

% ANNOTATION: Linde's comprehensive book on inflationary cosmology, including
% detailed treatments of phase transitions, bubble nucleation, and reheating.
% Essential context for the phase-transition interpretation of the Big Bang.
\bibitem{Linde1990}
A.~D.~Linde, \emph{Particle Physics and Inflationary Cosmology} (Harwood Academic Publishers, Chur, 1990). arXiv:hep-th/0503203.

% ANNOTATION: Coleman's original paper on vacuum decay in quantum field theory
% (without gravity). Establishes the instanton method for computing tunneling
% rates. Companion to the Coleman-De Luccia paper above.
\bibitem{Coleman1977}
S.~Coleman, ``Fate of the false vacuum: Semiclassical theory,'' \emph{Phys.\ Rev.\ D}\ \textbf{15}, 2929--2936 (1977).

% --- Induced gravity / emergent gravity ---

% ANNOTATION: Zee's paper on induced gravity, showing how gravitational dynamics
% can arise from quantum fluctuations of matter fields. Complements Sakharov's
% original induced gravity proposal and supports the mechanical-geometric linkage.
\bibitem{Zee1981}
A.~Zee, ``Spontaneously generated gravity,'' \emph{Phys.\ Rev.\ D}\ \textbf{23}, 858--866 (1981).

% ANNOTATION: Padmanabhan's paper on emergent gravity from the microscopic
% structure of spacetime. Develops the connection between thermodynamics and
% gravity, complementing Jacobson's approach.
\bibitem{Padmanabhan2010}
T.~Padmanabhan, ``Thermodynamical Aspects of Gravity: New insights,'' \emph{Rep.\ Prog.\ Phys.}\ \textbf{73}, 046901 (2010). arXiv:0911.5004.

% --- Additional analogue gravity experiments ---

% ANNOTATION: Lahav et al.'s realization of Unruh's proposal for detecting
% stimulated Hawking radiation in a water-wave analogue. Demonstrates that
% analogue gravity effects are observable in classical systems.
\bibitem{Lahav2010}
O.~Lahav \emph{et al.}, ``Realization of a Sonic Black Hole Analog in a Bose-Einstein Condensate,'' \emph{Phys.\ Rev.\ Lett.}\ \textbf{105}, 240401 (2010). arXiv:1008.0607.

% ANNOTATION: Weinfurtner et al.'s classical fluid experiment observing
% stimulated Hawking radiation in water waves. Complements BEC experiments.
\bibitem{Weinfurtner2011}
S.~Weinfurtner \emph{et al.}, ``Measurement of stimulated Hawking emission in an analogue system,'' \emph{Phys.\ Rev.\ Lett.}\ \textbf{106}, 021302 (2011). arXiv:1008.1911.

% ==============================================================================
% END OF BIBLIOGRAPHY ADDITIONS
% ==============================================================================

% ==============================================================================
% SUMMARY OF ADDITIONS
% ==============================================================================
%
% REQUIRED ENTRIES (7):
%   1. Volovik2003         - Universe in a Helium Droplet (canonical BEC cosmology)
%   2. Steinhauer2016      - Experimental Hawking radiation in BEC
%   3. ColemanDeLuccia1980 - Vacuum decay with gravity (nucleation theory)
%   4. GarayEtAl2000       - BEC analogue black holes (foundational proposal)
%   5. Israel1966          - Junction conditions (boundary physics foundation)
%   6. Jacobson1995        - Thermodynamics of spacetime (emergent gravity)
%   7. Volovik2008         - Fermi-point emergent physics
%
% ADDITIONAL RECOMMENDED ENTRIES (16):
%   8. Unruh1976           - Unruh effect (foundational)
%   9. Unruh1981           - Acoustic metric/dumb holes (foundational)
%  10. Visser1998          - Acoustic black holes (pedagogical)
%  11. DalfovoEtAl1999     - BEC theory review
%  12. Bogoliubov1947      - Bogoliubov excitations
%  13. Taub1948            - Relativistic Rankine-Hugoniot
%  14. Lichnerowicz1967    - Relativistic shock waves
%  15. Poulin2019          - Original EDE paper
%  16. SmithPoulinAmin2020 - Oscillating scalar field EDE
%  17. DESI2024            - Latest BAO results
%  18. DESICMB2024         - DESI+CMB Hubble tension analysis
%  19. Linde1990           - Inflationary cosmology / phase transitions
%  20. Coleman1977         - Vacuum decay (no gravity)
%  21. Zee1981             - Induced gravity
%  22. Padmanabhan2010     - Emergent gravity from thermodynamics
%  23. Lahav2010           - BEC analogue experiment
%  24. Weinfurtner2011     - Water-wave analogue experiment
%
% TOTAL: 24 new entries (7 required + 17 recommended)
% ==============================================================================

% ==============================================================================
% NOTES FOR INSERTION
% ==============================================================================
%
% Insert these entries into the existing bibliography in the following order:
%
% After "AnalogueGravityReview" (line 890), add:
%   - Volovik2003, Volovik2008 (emergent spacetime)
%   - Unruh1976, Unruh1981, Visser1998 (acoustic metrics)
%   - GarayEtAl2000, Steinhauer2016 (BEC analogues)
%   - Lahav2010, Weinfurtner2011 (additional experiments)
%   - Jacobson1995, Padmanabhan2010, Zee1981 (emergent/induced gravity)
%
% After "HillEDE2020" (line 825), add:
%   - Poulin2019, SmithPoulinAmin2020 (EDE foundations)
%
% After "eBOSS_BAO" (line 864), add:
%   - DESI2024, DESICMB2024 (latest BAO)
%
% New section "Phase transitions and nucleation" after BBN, add:
%   - ColemanDeLuccia1980, Coleman1977, Linde1990
%
% New section "Junction conditions and shock physics", add:
%   - Israel1966, Taub1948, Lichnerowicz1967
%
% New section "BEC theory" after analogue gravity, add:
%   - DalfovoEtAl1999, Bogoliubov1947
%
% ==============================================================================
